\section{\textsc{VesProof}: an embedded proof-script DSL}
\label{sec:dsl}

This section makes the recommendation of \cref{sec:design-embed}
concrete enough to evaluate: a tactic-oriented proof-script language
that lives \emph{inside} the Rust source as a block-shaped macro
adjacent to each \code{unsafe} site, its expansion into ordinary Rust
(ghost tokens the compiler itself checks), its deterministic
elaboration into Lean~4 statements, the diagnostics model, and three
worked examples from the real corpus --- including one where the proof
now fails \emph{at compile time}, as intended.

\subsection{Design lineage: the \code{asm!} pattern from \code{dotnetdll}}
\label{sec:dsl-lineage}

The project already contains a stylistic precedent for embedded DSLs
adjacent to Rust code. \dotnetrs{}'s metadata
layer is built on the author's \code{dotnetdll}
crate~\cite{dotnetdll}, whose \code{asm!} macro embeds a CIL assembler
directly in Rust:

\begin{lstlisting}[language=Rust]
Some(body::Method::new(asm! {
    load_string "Hello from dotnetdll!";
    call write_line;          // `write_line` is an ordinary Rust binding
    Return;
}))
\end{lstlisting}

Four properties of this macro family transfer directly to proof
scripts; \textsc{VesProof} adopts all four:

\begin{enumerate}[leftmargin=1.6em]
\item \textbf{Statement shape.} Each line is a keyword-led statement
  terminated by \code{;} --- not an attribute's nested meta-item tree.
  Scripts of ten steps stay readable; attribute syntax at that size
  does not.
\item \textbf{Native splicing.} Operands are ordinary Rust expressions
  resolved in the enclosing scope (\code{call write\_line}). No string
  quoting, no name mangling, ordinary hygiene.
\item \textbf{Binders into scope.} \code{asm!}'s \code{@label} form
  introduces Rust \code{let} bindings usable later in the block. Proof
  scripts need the same shape for premise witnesses:
  \code{have parked = ...;} binds a value later steps consume.
\item \textbf{One declarative source of truth.} \code{dotnetdll}'s
  entire \code{Instruction} enum --- variants, flag fields,
  constructors, pretty-printers --- is generated by a declarative macro
  from one compact specification. The \textsc{VesProof} vocabulary
  (goal and premise constructors) uses the same technique, generating the
  Rust token types \emph{and} the Lean declarations from a single
  versioned source, so the two sides cannot drift
  (\cref{sec:design-compart}).
\end{enumerate}

The earlier draft of this study proposed a
\code{\#[ves::proof(goal, using, lemma)]} attribute. That form is kept
only as a degenerate case (single-obligation sites with an
already-proved lemma); the primary surface is the script block below,
for the same reason \code{dotnetdll} did not encode CIL methods as
attribute arguments.

\subsection{Surface syntax}
\label{sec:dsl-syntax}

\code{ves::proof!} is an expression-position macro. Its body is a
proof script; its final statement is the protected expression, which
the macro re-emits \emph{verbatim} --- the block evaluates to whatever
the protected expression evaluates to.

\begin{lstlisting}[language=Rust,
  morekeywords={goal,have,show,then,trust,note,checked,infer}]
// SAFETY: prose stays; the script is the checked artifact.
let owner_id = ves::proof! {
    goal  read_frozen(ObjectInner::owner_id) at lock_ptr;
    have  parked   = stw::all_mutators_parked();
    have  alive    = arena::pinned_while_parked(parked, lock_ptr);
    have  wonce    = field::write_once(ObjectInner::owner_id);
    have  visible  = stw::handshake_publishes(parked, wonce);
    show  deref_frozen(alive, visible);
    then  unsafe { (*(*lock_ptr).as_ptr()).owner_id() }
};
\end{lstlisting}

The statement forms:

\begin{itemize}[leftmargin=1.4em]
\item \code{goal $g$ at $e_1, \dots$;} --- names the obligation, an
  instance of a vocabulary goal constructor applied to spliced Rust
  expressions and paths. \emph{Optional}: for the common shapes the
  vocabulary carries an inference table keyed on the protected
  expression's syntax (raw deref $\Rightarrow$ \code{read\_valid} /
  \code{write\_valid}; \code{from\_raw\_parts} $\Rightarrow$
  \code{slice\_valid}; \code{AtomicU$N$::from\_ptr} $\Rightarrow$
  \code{atomic\_access}), so short scripts start at \code{have}. An
  explicit \code{goal} always overrides inference, and the checker
  records which of the two produced the obligation.
\item \code{have $x$ = $t$;} --- acquires a premise witness. The
  right-hand side is a \emph{token expression}: a call into the
  vocabulary's ghost-token API (\cref{sec:dsl-tokens}), whose
  arguments may be earlier witnesses (the premise DAG is therefore
  type-checked by rustc) or spliced Rust values.
\item \code{trust $x$ = $\mathit{registry::name}$;} --- acquires a
  trusted-premise token. Compiles only if \code{name} is declared in
  \code{TRUSTED.toml} (\cref{sec:dsl-tokens}); counted by the ratchet.
\item \code{show $\mathit{tac}(x_1, \dots)$;} --- the discharge step: a
  \emph{tactic expression} naming a vocabulary tactic applied to
  witnesses. Tactics may nest (\code{show frame(alive, arith())}) and
  an escape hatch \code{show lean "..."} carries a raw Lean tactic
  block for sites that outgrow the surface vocabulary. \code{show} is
  a proof \emph{sketch}, not the proof: it elaborates to the head of a
  Lean tactic script (\cref{sec:dsl-elab}), and the site theorem in
  the proofs package may extend or replace it.
\item \code{then $e$} --- the protected expression or item, emitted
  verbatim. Exactly one per block, always last.
\item \code{note "..." ;} --- attached prose, carried into the manifest
  for reviewers; no semantic content.
\end{itemize}

An item-position twin, \code{ves::proof\_impl!\,\{ \dots\ then unsafe
impl \dots \}}, covers the 29 \code{unsafe impl} sites; a
\code{ves::exempt!\,(reason = "...")} form covers deliberate opt-outs.
The grammar is deliberately regular --- every statement is
\code{keyword \dots ;} --- so the parser is a few hundred lines of
\code{syn} in the shared \code{ves-syntax} crate, and \emph{the same
parser} is linked into both the proc macro and the external checker,
eliminating grammar skew by construction.

The four macros above are implemented (\code{ves-proof} commit
\code{85a6abc}, WP-7): \code{proof!} and \code{proof\_impl!} parse via
\code{ves\_syntax::parse\_proof\_script}/\code{parse\_proof\_impl\_script}
and expand \code{have}/\code{trust} statements to local bindings ---
\code{have $x$ = $t$;} becomes \code{let $x$ = ::ves\_tokens::$t$;}, and
\code{trust $x$ = registry::$c$;} becomes \code{let $x$ =
::ves\_tokens::registry::$c$(::ves\_tokens::grants::$C$);} where $C$ is
the upper-cased constructor name; \code{goal}, \code{show}, and
\code{note} are checker metadata and are omitted from the Rust
expansion entirely. The \code{then} tokens are spliced verbatim from
the parsed \code{ThenClause}, realizing the token-identity property
exactly as designed and confirmed by five expansion-snapshot fixtures.
\code{proof\_impl!} cannot introduce bindings at item scope, so its
premises type-check inside an anonymous \code{const \_: () = \{ fn
\_\_check() \{ \dots \} \};} block ahead of the verbatim \code{then}
item --- sound, but the isolated checker cannot capture caller-local
values or the protected item's own generics, a real (if narrow)
expressiveness limit worth remembering when scripting generic
\code{unsafe impl}s. \code{exempt!} validates its \code{reason}
argument and expands to nothing; \code{contract!} is, as planned, an
unchanged-input passthrough in this work package, with real
contract-token generation deferred to the contract-layer package
(\cref{sec:decomposition-wps}, item 10). Three compile-fail fixtures
exercise the premise-rejection paths this section relies on --- a
feature-gated checked constructor, a derived constructor with an
unmet prerequisite, and an unregistered trust path --- and all three
fail exactly the way \cref{sec:dsl-example-fail} predicts, as ordinary
\code{rustc} type errors rather than macro-specific diagnostics. One
naming detail is still open: the implemented crate exports bare
\code{proof}/\code{proof\_impl}/\code{exempt}/\code{contract} macros
(invoked as \code{ves\_macros::proof!}, etc.), not the \code{ves::}
facade used throughout this section's examples; assembling that facade
--- presumably a thin re-export module combining \code{ves-macros} and
\code{ves-tokens} --- is integration-bootstrap work
(\cref{sec:decomposition-wps}, item 9), not yet done.

\subsection{Premises as ghost tokens}
\label{sec:dsl-tokens}

The single largest improvement over the attribute draft is that
premises are not strings the external checker inspects --- they are
zero-sized Rust values with private constructors, generated from the
vocabulary source into the tiny runtime crate \code{ves-tokens}. A
\code{have} statement expands to an ordinary (inlined, empty,
release-free) function call returning a ZST:

\begin{lstlisting}[language=Rust]
// generated in ves-tokens (simplified)
pub struct Parked(());                      // ZST, private ctor
pub struct PinnedWhileParked<'a>(PhantomData<&'a ()>);

pub mod stw {
    #[inline(always)]
    pub fn all_mutators_parked() -> Parked {
        debug_assert!(crate::hooks::is_stw_in_progress());
        Parked(())
    }
}
\end{lstlisting}

Token constructors come in four trust classes, mirroring
\cref{sec:obligations-untrusted}:

\begin{itemize}[leftmargin=1.4em]
\item \emph{Static} --- acquirable anywhere; may carry a
  \code{debug\_assert!} shadow (as above) so the fixture suite
  dynamically spot-checks the premise for free.
\item \emph{Checked} --- the constructor is defined \emph{next to the
  runtime guard it names and under the same \code{cfg}}. This is the
  central rule: a guard that is compiled out cannot mint its
  token, so citing it is a \code{rustc} error in exactly the feature
  sets where the citation is false. Defect (a) of
  \cref{sec:system-defects} becomes unwritable
  (\cref{sec:dsl-example-fail}).
\item \emph{Trusted} --- the constructor's argument is a
  \code{TrustGrant} const that a build script generates from
  \code{TRUSTED.toml}; an unregistered trusted premise does not
  compile, and the ratchet counts grants, not comments.
\item \emph{Derived} --- constructors that consume other tokens
  (\code{pinned\_while\_parked(parked, ptr)}), so the dependency
  structure among premises is enforced by the type checker rather
  than by checker lint rules.
\end{itemize}

Two properties constrain this layer. \emph{Zero cost}: every token
is a ZST, every constructor is \code{\#[inline(always)]} and empty in
release; the pinned layout-size tests extend to assert
\code{size\_of} of every token is 0. \emph{Token-identity}: the macro
emits the \code{then}-expression's tokens unchanged --- \textsc{VesProof}
never rewrites the \code{unsafe} code it annotates --- and this is
testable by expansion snapshot (\code{cargo expand} diff against the
raw expression), which reduces the principal proc-macro risk to a CI
check.

The implemented \code{ves-tokens} (\code{ves-proof} commit
\code{73e6dea}) realizes the trust classes and both properties --- all
72 generated types carry a compile-time \code{size\_of == 0} assertion
kept exhaustive by a test that re-derives the type list from the
generated source, and every constructor is checked to be
\code{\#[inline(always)]} --- with three deliberate v0 simplifications
worth stating so that later work does not mistake them for settled
design. Static constructors carry no \code{debug\_assert!} shadow: the
predicates named above are \dotnetrs{} functions and the vocabulary has
no field to attach one, so the free dynamic spot-check is currently
forgone rather than refuted. No token is generic, so the
lifetime-branded form sketched above
(\code{PinnedWhileParked<'a>}) does not exist and F6 brands are not
carried in token types --- acceptable while the arena discipline is
enforced by \gcarena{}'s own branding, but it means a token cannot yet
witness \emph{which} arena it speaks about. Most consequentially,
checked constructors are gated by a Cargo feature of
\code{ves-tokens} rather than by the guard's own \code{cfg}, since the
two live in different crates. The guarantee is preserved only if
\dotnetrs{} forwards its guard \code{cfg} to that feature; if it
enables the feature unconditionally, a checked premise silently
degrades to a static one and the mechanism stops catching defect (a).
The integration bootstrap therefore owns the forwarding rule and a CI
leg that builds with the guard off and confirms the citing sites fail
to compile (\cref{sec:decomposition-wps}, item 9).

What remains for the \emph{external} structural checker is what rustc
cannot see: that every \code{unsafe} block in designated crates is
wrapped at all, that obligation hashes match the proof manifest
(\cref{sec:dsl-elab}), staleness policy, and the trusted-premise
budget. The external checker's role narrows while the guarantees
strengthen: name resolution, \code{cfg}-liveness, and premise
dataflow all move into the ordinary Rust build, where they are
checked in every feature configuration the CI matrix already covers.

\subsection{Binding, hygiene, spans, and \code{cfg}}
\label{sec:dsl-binding}

\paragraph{Binding by construction.} Because the script \emph{wraps}
the protected expression rather than sitting near it, there is no
adjacency convention to enforce and no way for refactors to separate
proof from code: move the expression and the script moves with it;
delete the expression and the script dies with it.

\paragraph{Hygiene.} \code{have}-bound witnesses are ordinary
\code{let} bindings introduced by the macro at the block's scope
(the \code{asm!} \code{@label} pattern); spliced expressions
(\code{lock\_ptr}, \code{self.layout}) resolve in the caller's scope
under standard \code{macro\_rules}/proc-macro hygiene. The macro
introduces no identifiers except the user's binders, and all
vocabulary paths are absolute (\code{::ves\_tokens::stw::...}), so
the expansion cannot capture or be captured.

\paragraph{Spans.} The proc macro preserves the token spans of every
script statement and of the protected expression, so (i) rustc errors
inside the \code{then}-expression point at the user's code exactly as
if the macro were absent, and (ii) the obligation manifest records a
real \code{file:line:col} span per statement, which the bridge uses to
map Lean-side failures back to source (\cref{sec:dsl-diag}).

\paragraph{\code{cfg}-awareness.} The expansion is subject to the same
\code{cfg} evaluation as the site itself, which resolves the
question the attribute draft handled with a bespoke checker rule:
a script under \code{\#[cfg(feature = "multithreading")]} simply does
not exist in single-threaded builds, and its obligations are extracted
per active feature set. The extraction pass runs once per CI matrix
leg (the matrix already exists), and the manifest keys obligations by
\code{(site, cfg-set)} --- the same site may legitimately carry
different proofs under different features, which is the reality of
this codebase (\code{ThreadSafeLock}'s two backends).

\subsection{Deterministic elaboration to Lean 4}
\label{sec:dsl-elab}

The pipeline from script to theorem has three mechanical stages and
one human one:

\begin{enumerate}[leftmargin=1.6em]
\item \textbf{Extraction.} \code{ves-check extract} parses the crate
  with \code{ves-syntax} (the same parser the macro links) and emits
  the \emph{obligation manifest}: for each site, the goal term (after
  inference), the premise list with trust classes, the tactic sketch,
  the span, the active \code{cfg} set, and two hashes --- the
  \emph{statement hash} over the canonicalized goal+premise terms, and
  the \emph{subject hash} over the protected expression's token
  stream. Canonicalization is total and documented (sorted premise
  sets, normalized paths, no source positions in the hash input), so
  extraction is deterministic across machines and macro-expansion
  order.
\item \textbf{Statement generation.} A Lean metaprogram in
  \code{VesCore} consumes the manifest and elaborates each goal into a
  \code{Prop} over the theory packages --- the same elaboration rules
  are versioned with the vocabulary, and the Lean side recomputes the
  statement hash from its elaborated term as a cross-check. For new
  obligations it emits a stub:
\begin{lstlisting}[language=Lean]
/-- site dotnet-value/src/object/mod.rs:288 (cfg: multithreading) -/
theorem site_9f3ac1 :
    ReadFrozen ObjectInner.owner_id lockPtr := by
  ves_sketch [all_mutators_parked, pinned_while_parked,
              write_once, handshake_publishes, deref_frozen]
  sorry
\end{lstlisting}
\item \textbf{Proof.} A human or agent replaces \code{sorry}, using
  the sketch tactics --- which are ordinary Lean tactic macros defined
  in the theory packages, i.e.\ \emph{untrusted automation whose
  output the Lean kernel checks}, the RefinedRust/Lithium
  discipline~\cite{refinedrust2024}. Reusable surface tactics
  (\code{deref\_frozen}, \code{frame}, \code{stw\_open}) are exactly
  reusable Lean lemma-combinators; growing the tactic library is
  proof engineering, never kernel engineering. This is where the
  bespoke-kernel temptation of \cref{sec:design-custom} is
  permanently discharged: the surface DSL gets arbitrarily
  domain-fitted while the TCB stays at Lean's kernel plus the
  elaboration rules.
\item \textbf{Publication.} \code{DotnetRsProofs}' CI exports the
  \emph{proof manifest}: obligation hash $\to$ theorem name,
  \code{sorry}-free flag, axioms reachable from the proof
  (\code{\#print axioms} transitively), toolchain and vocabulary
  versions. \dotnetrs{}'s CI diffs its freshly extracted obligation
  manifest against this artifact --- no Lean toolchain anywhere in the
  VM's build (\cref{sec:design-compart}).
\end{enumerate}

Statement-hash discipline gives the drift semantics: change a
script's goal or premises and the hash moves --- the site reverts to
\emph{unproved} (blocking or ratcheted, per policy) until the Lean
side catches up; change only the protected expression and the subject
hash moves --- a weaker ``re-audit'' state, since the statement may
still hold, flagged rather than blocked by default. Both states are
visible in the checker's report, which is the mechanized version of
the re-audit triggers \cref{sec:risks-fidelity} asks for.

\subsection{Diagnostics}
\label{sec:dsl-diag}

Three failure surfaces, three reporting paths, all span-carrying:

\begin{itemize}[leftmargin=1.4em]
\item \emph{Script errors} (unknown vocabulary item, malformed
  statement, missing \code{then}) are proc-macro parse errors ---
  ordinary rustc diagnostics at the offending token, with the
  vocabulary crate supplying ``did you mean'' candidates.
\item \emph{Premise errors} (token constructor cfg'd out, unregistered
  trust grant, dataflow type mismatch) are ordinary type errors at the
  \code{have} statement's span --- rustc's own machinery, in every
  matrix leg.
\item \emph{Proof-state errors} (hash mismatch, \code{sorry},
  staleness, budget breach) come from \code{ves-check}, which formats
  them rustc-style against the manifest's recorded spans, including
  the Lean theorem name and the last-good hash so the fix is
  mechanical to locate.
\end{itemize}

\subsection{Worked example 1: the cross-arena rebrand (F1+F6)}
\label{sec:dsl-example-rebrand}

The principal accepted-risk boundary (\code{heap.rs:436--445}) in
script form:

\begin{lstlisting}[language=Rust,
  morekeywords={goal,have,show,then,trust,note}]
ves::proof! {
    goal  rebrand_confined(ptr => tracing_arena) in trace_call;
    have  parked = stw::all_mutators_parked();
    have  pinned = stw::arenas_pinned_while_parked(parked);
    have  fresh  = arena::recorded_cross_ref(ptr, generation);
    have  scoped = scope::no_escape_of_result();
    show  rebrand_ok(pinned, fresh, scoped);
    then  unsafe { Gc::from_ptr(ptr.as_ptr()).trace(cc); }
}
\end{lstlisting}

\code{scope::no\_escape\_of\_result()} is the one premise discharged
syntactically by the macro itself (the handle is consumed within the
\code{then}-expression; no binder escapes), recorded in the manifest
as macro-checked. The interesting premise is
\code{arenas\_pinned\_while\_parked} --- a derived token whose Lean
counterpart is \emph{the} protocol theorem
(\cref{sec:risks-concurrency}): in phase \code{parked}, no arena
teardown transition exists and \code{unregister\_arena} blocks on
\code{active\_leases}. Until that theorem lands, the site compiles
with \code{trust pinned = registry::stw\_pins\_arenas;} instead ---
one changed line, one ratchet entry, and the flip from
\code{trusted} to \code{proved} later is a one-line diff whose hash
the checker verifies.

\subsection{Worked example 2: tracing completeness (F5+F2)}

The item-position form on \code{FieldStorage}'s \code{Collect} impl:

\begin{lstlisting}[language=Rust,
  morekeywords={goal,have,show,then,trust,note}]
ves::proof_impl! {
    goal  trace_complete(FieldStorage) via gc_desc(layout);
    have  parked   = stw::all_mutators_parked();
    have  visible  = stw::handshake_publishes(parked,
                        field::all_of(FieldStorage));
    trust faithful = registry::gc_desc_faithful;
    show  trace_via_desc(visible, faithful);
    then  unsafe impl<'gc> Collect<'gc> for FieldStorage { /* ... */ }
}
\end{lstlisting}

Elaboration splits the goal into the per-type completeness statement
(every reference the bytes contain sits at an offset \code{GcDesc}
marks) and the data-race side condition (discharged by the same STW
theorem as example 1 plus the handshake happens-before lemma, closing
the gap the prose hand-waves in F7). \code{gc\_desc\_faithful} starts
life as a trust grant and is retired in Phase 4 by the
\code{GcDescFaithful} theorem against the mechanized \ecma{II.10.7}
rules --- provable once per layout algorithm, and testable against the
executable model on the fixture corpus's types \emph{before} it is
proved.

\subsection{Worked example 3: a proof that fails at compile time (F4)}
\label{sec:dsl-example-fail}

The defective site from \cref{sec:system-defects}(a), scripted
honestly:

\begin{lstlisting}[language=Rust,
  morekeywords={goal,have,show,then,trust,note}]
let v = ves::proof! {
    // goal inferred: atomic_access(u32) at ptr
    have aligned = checked::alignment_guard(ptr, align_of::<u32>());
    //             ^^^^^^^^^^^^^^^^^^^^^^^ rustc E0425 in default
    //             features: constructor exists only under
    //             cfg(feature = "memory-validation"), same cfg as
    //             the guard it witnesses.
    show atomic_from_witness(aligned);
    then unsafe { AtomicU32::from_ptr(ptr as *mut u32) }.load(ord)
};
\end{lstlisting}

In the attribute draft this was caught by checker rule 3; here it is
not writable at all --- the token constructor is \code{cfg}-mirrored
with \code{validate\_alignment}, so the build that ships without the
guard is the build in which this script does not compile. The
author's options are unchanged and all sound: acquire the witness
from the always-on guard (\code{is\_ptr\_aligned\_to\_field} returns
its token in every configuration), push a contract token onto the
\code{pub unsafe fn} boundary so every caller supplies it, or route
through the internally-checked \code{Atomic::load\_field}. The
earliest possible detection moved from CI to the editor.

\subsection{Treatment of template comments}

Unchanged in substance from the attribute draft, upgraded in
mechanism: the 342 boilerplate comments collapse into contract tokens
on the $\sim$62 documented \code{unsafe fn} --- the function's
\code{\#\,Safety} section becomes a \code{ves::contract!} declaration
minting a per-function contract token type; each caller's script
acquires it (\code{have c = contract::write\_to\_unmanaged(bounds);})
by supplying the bounds-check witness; the callee's own script
consumes \code{c} as its ambient precondition. One contract, one
lemma, dozens of sites, and the chain from caller check to callee
deref is a typed dataflow rustc verifies --- the ``documented
preconditions and preceding bounds/origin checks'' template sentence
becomes a value with a definition. The bimodality of the corpus
(\cref{sec:system-corpus}) remains the reason the campaign is
tractable: $\sim$25 goal constructors, $\sim$40--60 real lemmas, a
handful of protocol theorems.
